% AAAI-27 Student Abstract and Poster Program draft
% Track deadline: 28 September 2026
% Requirement: maximum 2 pages INCLUDING references, AAAI two-column camera-ready style.
% IMPORTANT before submission:
%   1. Compile with the official AAAI-27 author kit (aaai27.sty).
%   2. Replace the phone-number placeholder below.
%   3. Check the final PDF is no more than 2 pages including references.
%   4. Submit a separate supplemental PDF/ZIP (<=20 MB), as requested by the CFP.

\documentclass[letterpaper]{article}

% Use the official AAAI-27 Author Kit files: aaai2027.sty and aaai2027.bst.
\usepackage{aaai2027}

\usepackage[hyphens]{url}
\usepackage{graphicx}
\usepackage{booktabs}
\usepackage{microtype}

\urlstyle{rm}
\def\UrlFont{\rm}

% AAAI papers normally use unnumbered sections.
\setcounter{secnumdepth}{0}

\title{Evaluating Visual Faithfulness in Medical Vision--Language Models}

\author{
    Dawud Izza\textsuperscript{\rm 1},
    Frank Soboczenski\textsuperscript{\rm 1}
}
\affiliations{
    \textsuperscript{\rm 1}University of York, York, United Kingdom\\
    Postal address: University of York, York, YO10 5GH, United Kingdom\\
    Primary author phone: [ADD PHONE NUMBER]\\
    dawud.izza@york.ac.uk, frank.soboczenski@york.ac.uk
    % Optional: add a personal/research URL for Dawud if desired.
}

\begin{document}

\maketitle

\begin{abstract}
Medical vision--language models (VLMs) can answer questions about clinical images and generate explanations for their predictions, but plausible outputs do not necessarily establish that a model is using the relevant visual evidence. This work investigates visual faithfulness in medical VLMs through three complementary questions: whether answers depend on the supplied image, whether explanations identify evidence that influences model behaviour, and whether confidence responds appropriately when visual evidence is reduced or altered. Gastrointestinal (GI) endoscopy visual question answering (VQA) provides the case study. The proposed evaluation uses controlled image interventions---including correct, shuffled, and neutral visual inputs---together with targeted perturbation of explanation-relevant regions and confidence analysis. An initial implementation provides a resource-efficient VLM fine-tuning and inference pipeline and components for visual attribution and confidence extraction; empirical evaluation is the next stage of the work. The intended contribution is a reproducible evaluation methodology for distinguishing plausible explanations from explanations that are behaviourally linked to clinically relevant visual evidence.
\end{abstract}

\section{Motivation and Research Gap}

GI endoscopy is a visually demanding clinical setting in which relevant findings can vary substantially with anatomy, viewpoint, illumination, image quality, bowel preparation, and procedural artefacts~\cite{beg2017quality}. Recent datasets and shared tasks such as Kvasir-VQA, Kvasir-VQA-x1, and Medico provide image--question--answer pairs for evaluating multimodal reasoning over gastrointestinal imagery~\cite{gautam2024kvasirvqa,gautam2025kvasirvqa,gautam2025medico}. Parameter-efficient adaptation has also shown promise for GI VQA~\cite{peter2026parameterefficient}.

However, answer quality alone is insufficient to establish that a VLM is using the image in the way its explanation implies. Medical VLM evaluations have identified hallucination and weak visual grounding as continuing concerns~\cite{chang2025medheval,liu2026medical}. More generally, a model may exploit linguistic or dataset regularities, continue answering when visual evidence is absent, or provide an apparently coherent explanation whose highlighted evidence is not important to the prediction. In a clinical setting, such explanation--evidence mismatch is especially important because a plausible explanation may encourage reliance without demonstrating visual faithfulness.

This work therefore asks three linked questions:
\begin{enumerate}
    \item \textbf{Visual dependence:} Does the answer depend on the supplied image?
    \item \textbf{Explanation faithfulness:} Does the explanation identify evidence that influences the answer?
    \item \textbf{Confidence sensitivity:} Does model confidence respond when relevant visual evidence is reduced or altered?
\end{enumerate}

\section{Evaluation Approach}

The study uses GI VQA as a controlled case study. A fixed question is evaluated under multiple visual-evidence conditions. With the \emph{correct image}, task-relevant evidence is available; with a \emph{shuffled image}, the question is paired with different but plausible GI imagery; and with a \emph{neutral image}, task-relevant visual information is largely removed. Comparing answer quality and confidence across these conditions provides a behavioural test of whether predictions depend on the intended visual input rather than primarily on language cues.

\begin{table}[t]
\centering
\small
\caption{Core visual-evidence interventions.}
\label{tab:conditions}
\begin{tabular}{p{0.24\columnwidth} p{0.68\columnwidth}}
\toprule
\textbf{Condition} & \textbf{Purpose} \\
\midrule
Correct image & Preserve task-relevant visual evidence. \\
Shuffled image & Replace the paired image with different plausible GI content. \\
Neutral image & Remove most task-relevant visual information while retaining an image input. \\
\bottomrule
\end{tabular}
\end{table}

Explanation faithfulness is evaluated separately from answer correctness. Visual regions identified as important by an attribution or explanation method are selectively altered, for example through masking or deletion. The analysis then measures whether the answer or its confidence changes. If an explanation identifies evidence that is genuinely influential, perturbing that evidence should produce a measurable behavioural effect. This moves beyond assessing whether an explanation appears clinically plausible.

Finally, confidence is treated as an additional behavioural signal. Prior work has shown that predictive performance and uncertainty need not be well aligned in VLMs~\cite{kostumov2024uncertainty}. The study therefore examines whether confidence decreases when relevant visual evidence is weakened, removed, or replaced, and whether this sensitivity is associated with changes in answer correctness and explanation behaviour.

\section{Current Status and Intended Contribution}

Preliminary exploratory work examined GI polyp segmentation using convolutional and transformer-based models together with saliency-based visual explanations. This established an initial image-analysis pipeline and motivated the move from pixel-level predictive performance toward questions of evidence grounding and explanation faithfulness.

Current development focuses on a GI VQA baseline using Kvasir-VQA-x1. The implemented architecture includes a resource-efficient VLM fine-tuning and inference pipeline together with components for extracting visual attribution and answer-confidence signals. The controlled faithfulness experiments described above have not yet been completed; empirical evaluation is the next stage of the work.

The intended contribution is a reproducible methodology for evaluating whether medical VLM behaviour is visually grounded across answers, explanations, and confidence estimates. GI endoscopy provides a concrete safety-critical testbed, while the evaluation principles are intended to generalise to other medical multimodal tasks. A later validation stage will examine whether the proposed measures align with clinician judgement, and future work may extend the same principles to multi-step and agentic medical AI systems.

\bibliography{references}

\end{document}
